\section{Singular denominators}\label{app:C.4}

An energy denominator $1/(E_i - E_m + \ii0^+)$ is the transform of a causal time evolution, \cref{eq:C.1-causal}, and it is not an ordinary function at $E_m = E_i$. \Cref{sec:2.4} split it into a principal value and a delta function, \cref{eq:2.4-SP}, by taking the limit of a finite width. It also gave the squared denominator a meaning through the Hadamard finite part, \cref{eq:2.4-fp}, with the expansion \cref{eq:2.4-split} of a Lorentzian in powers of its width. Here we prove the splitting for any smooth test function, collect the rules for manipulating these distributions, and give the correction terms of the Lorentzian limits, which control how a broadened spectrum approaches its zero-width limit.

\subsection{The Sokhotski--Plemelj formula}

Let $w(x)$ be a test function, bounded, integrable and differentiable at $x = 0$; a Lipschitz condition $|w(x) - w(0)| \le c|x|$ near the origin is enough. The principal value is the symmetric limit
\begin{equation}
  \operatorname{PV}\!\int\dd x\,\frac{w(x)}{x} = \lim_{x_{\mathrm{c}}\to0}\int_{|x| > x_{\mathrm{c}}}\dd x\,\frac{w(x)}{x} = \int\dd x\,\frac{w(x) - w(0)\,\Theta(X_{\mathrm{c}} - |x|)}{x} ,
  \label{eq:C.4-PV}
\end{equation}
where the second form holds for any fixed $X_{\mathrm{c}} > 0$. On the symmetric interval $|x| < X_{\mathrm{c}}$ the subtracted term integrates to zero because $1/x$ is odd, and the remaining integrand is bounded near the origin by the Lipschitz condition. The second form shows that the principal value is finite. Folding the negative half line onto the positive one gives $\operatorname{PV}\!\int\dd x\,w(x)/x = \int_0^\infty\dd x\,[w(x) - w(-x)]/x$, whose integrand is finite at the origin: a principal-value integral over a tabulated spectrum is computed by folding, not by cutting out an interval around the pole.

For a finite width $\varGamma_{\mathrm{w}} > 0$, the real and imaginary parts of $1/(x + \ii\varGamma_{\mathrm{w}}/2)$ act on $w$ separately. The imaginary part is $-\pi$ times the normalized Lorentzian $(\varGamma_{\mathrm{w}}/2\pi)/(x^2 + \varGamma_{\mathrm{w}}^2/4)$, which is the squeezed family $t\,f(tx)$ of \cref{app:C.1} with $f(u) = \pi^{-1}/(1 + u^2)$ and $t = 2/\varGamma_{\mathrm{w}}$. Its integral against $w$ tends to $w(0)$ by the dominated-convergence argument given there. The real part is $x/(x^2 + \varGamma_{\mathrm{w}}^2/4)$. Subtracting $w(0)\,\Theta(X_{\mathrm{c}} - |x|)$, which integrates to zero against this odd kernel,
\begin{equation*}
  \int\dd x\,\frac{x\,w(x)}{x^2 + \varGamma_{\mathrm{w}}^2/4} = \int\dd x\,\frac{x^2}{x^2 + \varGamma_{\mathrm{w}}^2/4}\;\frac{w(x) - w(0)\,\Theta(X_{\mathrm{c}} - |x|)}{x} .
\end{equation*}
The first factor lies between 0 and 1 and tends to 1 at every $x \neq 0$; the second is integrable and independent of $\varGamma_{\mathrm{w}}$. By dominated convergence the integral tends to the second form of \cref{eq:C.4-PV}. Together,
\begin{equation}
  \lim_{\varGamma_{\mathrm{w}}\to0^+}\int\dd x\,\frac{w(x)}{x - x_0 \pm \ii\varGamma_{\mathrm{w}}/2} = \operatorname{PV}\!\int\dd x\,\frac{w(x)}{x - x_0} \mp \ii\pi\,w(x_0) ,
  \label{eq:C.4-SP}
\end{equation}
after a shift of the variable by $x_0$. With the lower sign this is the complex conjugate statement. Written as distributions, $1/(x \pm \ii0^+) = \operatorname{PV}(1/x) \mp \ii\pi\delta(x)$, the Sokhotski--Plemelj formula, \cref{eq:2.4-SP}. The proof used only that $w$ is integrable and Lipschitz at $x_0$; a sum over a dense set of final states, as in the level shift of \cref{eq:2.4-WW}, satisfies it once the levels are replaced by a smooth density of states (\cref{app:C.2}).

The formula has an equivalent form for functions of a complex variable. Let
\begin{equation}
  \mathcal{F}(z) = \frac{1}{2\pi\ii}\int\dd x\,\frac{w(x)}{x - z} ,
  \label{eq:C.4-cauchy}
\end{equation}
a Cauchy integral, which is analytic in $z$ above and below the real axis. Approaching a point $x_0$ of the axis from above, $z = x_0 + \ii0^+$, and from below, \cref{eq:C.4-SP} gives
\begin{equation}
  \mathcal{F}(x_0 + \ii0^+) - \mathcal{F}(x_0 - \ii0^+) = w(x_0), \qquad
  \mathcal{F}(x_0 + \ii0^+) + \mathcal{F}(x_0 - \ii0^+) = \frac{1}{\pi\ii}\operatorname{PV}\!\int\dd x\,\frac{w(x)}{x - x_0} .
  \label{eq:C.4-jump}
\end{equation}
The jump across the axis is the density $w$, and the average of the two sides is its principal-value transform. Every response function of the book has this structure: its imaginary part on the real axis is the density of transitions, and its real part is a principal-value integral over that density. \Cref{app:C.5} draws the consequences, the Kramers--Kronig relations.

\subsection{Rules for the distributions}

Since a distribution is defined by its action on test functions, identities among them are proved by integrating against $w$.

The product of $\delta(x)$ with a smooth function $h(x)$ is $h(0)\delta(x)$; in particular $x\,\delta(x) = 0$ and $x\operatorname{PV}(1/x) = 1$, since $\operatorname{PV}\!\int\dd x\,x\,w(x)/x = \int\dd x\,w(x)$. The difference of the two boundary values is a delta function,
\begin{equation}
  \frac{1}{x - \ii0^+} - \frac{1}{x + \ii0^+} = 2\pi\ii\,\delta(x) ,
  \label{eq:C.4-difference}
\end{equation}
which is \cref{eq:C.4-SP} with both signs. Derivatives act by moving them onto the test function, $\int\dd x\,T'(x)\,w(x) = -\int\dd x\,T(x)\,w'(x)$, the definition that makes integration by parts valid for distributions~\cite[Eq.~(1.16.8)]{olver2010nist}. Hence $\dd\ln|x|/\dd x = \operatorname{PV}(1/x)$, as follows from $\int\dd x\,\ln|x|\,w'(x)$ split at $\pm x_{\mathrm{c}}$, where the boundary terms $\ln x_{\mathrm{c}}\,[w(x_{\mathrm{c}}) - w(-x_{\mathrm{c}})]$ vanish as $x_{\mathrm{c}}\to0$.

The derivative of the principal value is minus the finite part. Splitting at $\pm x_{\mathrm{c}}$ and integrating by parts,
\begin{equation*}
  -\int_{|x| > x_{\mathrm{c}}}\dd x\,\frac{w'(x)}{x} = \frac{w(x_{\mathrm{c}}) + w(-x_{\mathrm{c}})}{x_{\mathrm{c}}} - \int_{|x| > x_{\mathrm{c}}}\dd x\,\frac{w(x)}{x^2} ,
\end{equation*}
and as $x_{\mathrm{c}}\to0$ the right side becomes minus the bracket of \cref{eq:2.4-fp}, because $w(x_{\mathrm{c}}) + w(-x_{\mathrm{c}}) = 2w(0) + O(x_{\mathrm{c}}^2)$. Therefore
\begin{equation}
  \frac{\dd}{\dd x}\operatorname{PV}\frac1x = -\operatorname{fp}\frac{1}{x^2} .
  \label{eq:C.4-dPV}
\end{equation}
The finite part is thus the derivative of a distribution we already have, and it obeys the usual rules: $x^2\operatorname{fp}(1/x^2) = 1$, $x\operatorname{fp}(1/x^2) = \operatorname{PV}(1/x)$, and $\operatorname{fp}\!\int\dd x\,w(x)/x^2 = 0$ for constant $w$, the last one by the explicit form of \cref{eq:2.4-fp}.

The square of a denominator follows from the derivative. Since $1/(x + \ii\varGamma_{\mathrm{w}}/2)^2 = -\dd[1/(x + \ii\varGamma_{\mathrm{w}}/2)]/\dd x$ at every finite width, and derivatives commute with the limit when both act on a test function, \cref{eq:2.4-SP,eq:C.4-dPV} give
\begin{equation}
  \frac{1}{(x + \ii0^+)^2} = \operatorname{fp}\frac{1}{x^2} + \ii\pi\,\delta'(x) .
  \label{eq:C.4-square}
\end{equation}
A double pole on one side of the axis is a well-defined distribution. The modulus squared $|x + \ii\varGamma_{\mathrm{w}}/2|^{-2}$ is a different object, because it is the product of a denominator and its conjugate, whose poles lie on opposite sides of the axis and pinch it as $\varGamma_{\mathrm{w}}\to0$. Its limit does not exist; its expansion in $\varGamma_{\mathrm{w}}$ starts with $2\pi\delta(x)/\varGamma_{\mathrm{w}}$, \cref{eq:2.4-split}. Products of distributions are in general undefined, and the squared delta function of a rate always comes from such a pinch.

\subsection{Corrections to the Lorentzian limit}

The normalized Lorentzian of full width $\varGamma_{\mathrm{w}}$ at half maximum is
\begin{equation}
  \delta_{\varGamma}(x) = \frac{1}{\pi}\,\frac{\varGamma_{\mathrm{w}}/2}{x^2 + \varGamma_{\mathrm{w}}^2/4} = -\frac1\pi\Imag\frac{1}{x + \ii\varGamma_{\mathrm{w}}/2} ,
  \label{eq:C.4-lorentz}
\end{equation}
and \cref{eq:2.4-split} states how it approaches the delta function. Multiplied by $\varGamma_{\mathrm{w}}/2\pi$, it reads
\begin{equation}
  \int\dd x\,\delta_{\varGamma}(x)\,w(x) = w(0) + \frac{\varGamma_{\mathrm{w}}}{2\pi}\operatorname{fp}\!\int\dd x\,\frac{w(x)}{x^2} - \frac{\varGamma_{\mathrm{w}}^2}{8}\,w''(0) + O(\varGamma_{\mathrm{w}}^3) .
  \label{eq:C.4-lorentzcorr}
\end{equation}
The first correction is linear in the width and involves the whole of $w$, through the finite part. It is the price of the slow $1/x^2$ tails, which give weight $\varGamma_{\mathrm{w}}/(\pi|x|)$ to the region beyond any $|x|$. A Gaussian of the same width has no linear term: its first correction is proportional to $\varGamma_{\mathrm{w}}^2\,w''(0)$, and it depends on $w$ only near the origin (\cref{app:C.1}). For $w(x) = 1/(1 + x^2)$, the Lorentzian integral is exactly $1/(1 + \varGamma_{\mathrm{w}}/2)$, by closing the contour in either half plane. The finite part of $\int\dd x\,w(x)/x^2$ is $-\pi$, since $1/[x^2(1 + x^2)] = 1/x^2 - 1/(1 + x^2)$ and the finite part of $\int\dd x/x^2$ vanishes, and $w''(0) = -2$. \Cref{eq:C.4-lorentzcorr} then reads $1 - \varGamma_{\mathrm{w}}/2 + \varGamma_{\mathrm{w}}^2/4 + \dots$, the expansion of the exact result. At $\varGamma_{\mathrm{w}} = 0.2$ the linear term is a \SI{10}{\percent} correction, and a Gaussian of the same full width at half maximum, whose correction is $\tfrac12w''(0)$ times its variance $(\varGamma_{\mathrm{w}}/2.355)^2$, would make an error of \SI{0.7}{\percent}.

The Lorentzian is the line shape of exponential decay. A state that decays as $\ee^{-\varGamma_{\mathrm{w}}|t|/2}$ in amplitude has, by \cref{eq:A.3-time}, the spectrum $\int\dd t\,\ee^{\ii xt}\ee^{-\varGamma_{\mathrm{w}}|t|/2} = 2\pi\delta_{\varGamma}(x)$. The convolution theorem, \cref{eq:C.1-conv} in its time form, then fixes the shape of combined lines. Two independent decay processes multiply in time and convolve in energy, $\ee^{-\varGamma_1|t|/2}\ee^{-\varGamma_2|t|/2} = \ee^{-(\varGamma_1 + \varGamma_2)|t|/2}$. The convolution of Lorentzians is a Lorentzian whose width is the sum of the two, $\delta_{\varGamma_1}\ast\delta_{\varGamma_2} = \delta_{\varGamma_1 + \varGamma_2}$; widths of independent channels add, as \cref{sec:2.4} assumed for the total width of an intermediate state. A Lorentzian convolved with a Gaussian, the Voigt profile, keeps the $1/x^2$ tails of the Lorentzian, so the linear correction of \cref{eq:C.4-lorentzcorr} survives any additional Gaussian broadening. A width that depends on energy, $\varGamma_{\mathrm{w}}(x)$, keeps the delta limit of \cref{eq:C.4-SP} at its value at $x_0$, provided it varies little over a range $\varGamma_{\mathrm{w}}$ around the pole, and its slope enters the finite-width corrections.

